\chapter{Dysarthria (Slurred Speech)}

\section*{Definition}
Dysarthria is a motor speech disorder characterized by poor articulation, slow or slurred speech, and reduced intelligibility due to weakness, incoordination, or paralysis of the muscles involved in speech production. It results from neurological damage affecting the central or peripheral nervous system.

\section*{What Happens in Your Body}
Dysarthria arises from damage to any component of the motor speech system: upper motor neuron (corticobulbar tract), lower motor neuron (cranial nerves V, VII, IX, X, XII), cerebellum, basal ganglia, or neuromuscular junction. Lesion location determines dysarthria type: spastic (UMN), flaccid (LMN), ataxic (cerebellar), hypokinetic (basal ganglia), hyperkinetic (dystonia/chorea), or mixed.

\section*{Causes}
\begin{itemize}
  \item Stroke (ischemic or hemorrhagic -- most common acute cause)
  \item Traumatic brain injury
  \item Parkinson disease (hypokinetic dysarthria)
  \item Multiple sclerosis
  \item Amyotrophic lateral sclerosis (ALS -- mixed dysarthria)
  \item Brain tumor (especially affecting brainstem, cerebellum, or motor cortex)
  \item Huntington disease (hyperkinetic dysarthria)
  \item Cerebellar degeneration (ataxic dysarthria -- alcohol, paraneoplastic)
  \item Myasthenia gravis (fluctuating dysarthria with fatigue)
  \item Guillain-Barre syndrome (flaccid dysarthria)
  \item Cerebral palsy (developmental dysarthria)
  \item Medications: Sedatives, anticonvulsants, antipsychotics
  \item Alcohol intoxication (acute, reversible)
\end{itemize}

\section*{When to See a Doctor}
See your doctor if:
\begin{itemize}
  \item Acute onset of dysarthria (suspect stroke or TIA -- emergency)
  \item Dysarthria with sudden facial droop, arm weakness, or confusion (stroke symptoms)
  \item Progressive worsening of speech over time
  \item Dysarthria with difficulty swallowing (dysphagia) or choking
  \item Fluctuating dysarthria that worsens with fatigue (myasthenia gravis)
  \item Associated limb weakness, numbness, or visual changes (multiple sclerosis)
\end{itemize}

\section*{Related Symptoms}
\begin{itemize}
  \item Aphasia (language difficulty)
  \item Dysphagia (difficulty swallowing)
  \item Facial weakness
  \item Limb weakness or numbness
  \item Ataxia (loss of coordination)
  \item Tremor
  \item Cognitive changes
  \item Drooling
  \item Emotional lability
\end{itemize}

\section*{Self-Care / Home Management}
\begin{itemize}
  \item For acute onset, call emergency services immediately (FAST: Face, Arms, Speech, Time).
  \item Speak slowly and use gestures or writing if speech is unintelligible.
  \item Maintain eye contact and allow extra time for communication.
  \item Reduce background noise to improve communication.
  \item Seek speech-language pathology evaluation for non-acute cases.
  \item Do not self-treat; underlying cause requires medical diagnosis.
\end{itemize}

\section*{How Your Doctor Will Evaluate You}
\begin{itemize}
  \item Neurological examination assessing cranial nerves, motor function, coordination, and gait.
  \item Speech evaluation by speech-language pathologist (intelligibility, rate, resonance, articulation).
  \item MRI brain (preferred) or CT head for stroke, tumor, or demyelinating disease.
  \item Electromyography (EMG) and nerve conduction studies for motor neuron disease.
  \item Acetylcholine receptor antibodies for myasthenia gravis.
  \item Genetic testing for hereditary disorders.
\end{itemize}

\section*{Treatment Options}
\begin{itemize}
  \item Treat underlying cause: Thrombolysis for ischemic stroke, carbidopa-levodopa for Parkinson disease, immunotherapy for multiple sclerosis.
  \item Speech-language therapy focused on articulation exercises, breathing support, and pacing techniques.
  \item Augmentative and alternative communication (AAC) devices for severe cases.
  \item Pharmacological management for specific etiologies.
  \item Supportive care for progressive conditions (ALS, Huntington disease).
\end{itemize}

\clearpage
